\documentclass[pdflatex,sn-mathphys-num]{sn-jnl}

\usepackage{graphicx}%
\usepackage{multirow}%
\usepackage{amsmath,amssymb,amsfonts}%
\usepackage{amsthm}%
\usepackage{mathrsfs}%
\usepackage[title]{appendix}%
\usepackage{xcolor}%
\usepackage{textcomp}%
\usepackage{manyfoot}%
\usepackage{booktabs}%
\usepackage{algorithm}%
\usepackage{algorithmicx}%
\usepackage{algpseudocode}%
\usepackage{listings}%

\raggedbottom

\begin{document}

\title[Article Title]{Explainable and Adaptable IoT Automation Generation for Personalized Clinical Care}

\author*[1]{\fnm{James} \sur{Nicholls}}\email{jnich6@uwo.ca}

\author*[1,2]{\fnm{Roy} \sur{Eagleson}}\email{eagleson@uwo.ca}

\author*[2,3]{\fnm{Sandrine} \sur{De Ribaupierre}}\email{sderibau@uwo.ca}

\affil[1]{\orgdiv{Faculty of Engineering, Department of AI and Software Engineering}, \orgname{Western University}, \orgaddress{\street{1151 Richmond Street}, \city{London}, \postcode{N6A3K7}, \state{Ontario}, \country{Canada}}}

\affil[2]{\orgdiv{Center for Brain and Mind}, \orgname{Western University}, \orgaddress{\street{1151 Richmond Street}, \city{London}, \postcode{N6A3K7}, \state{Ontario}, \country{Canada}}}

\affil[2]{\orgdiv{Clinical Neurological Sciences}, \orgname{Western University}, \orgaddress{\street{1151 Richmond Street}, \city{London}, \postcode{N6A3K7}, \state{Ontario}, \country{Canada}}}

\abstract{As daily life becomes increasingly digital, there is a growing need to automate fragmented Internet of Things (IoT) infrastructure. Many tools can be used to create Personal Life Automations (PLAs), automating actions across these services. These tools are time-consuming and costly to implement, since PLAs must be created by their end-users to achieve personalization. Users require domain and technical knowledge of the systems they are automating. We introduce a ML pipeline to adapt and automate the changing patterns in IoT systems through generating automation code. This metholodology will be applied to clinical care orchestration, implementing a perioperative workflow commonly cited as ”OR 2.0”. Our novel approach enables distributed learning while privatizing training data through Federated Learning (FL) with Differential Privacy (DP). Human feedback is incorporated so generalized models can rapidly adapt for individual users. Our system, while using the Smart Home as a foundation, will be applicable to any IoT environment or domain. We chose to build upon the Home Assistant platform for its broad device compatibility enabled by open-source contributions. We demonstrate performance even with on-premises budget inference hardware.}

\keywords{Human AI Interaction, AI Safety, Artificial Intelligence, Smart Home, Internet of Things, Reinforcement Learning, Multi-Layer Perceptron, Decision Tree, Deep Q-Network, Large Language Model, Operating Room 2.0, Automation, Personal Life Automation, Explainable AI, Direct Preference Optimization}

\maketitle

\section{Introduction}\label{}
Your mom \cite{abadi_deep_2016}

\bibliography{jnich6-aissh-2026-pub}

\end{document}